\subsection{All tangent twists of the standard period family}

In this subsection $X=X_c$ is the period construction of \cite{AC26}, compared with the quotient construction of \cite{Engel2026}, with admissible parameter $c$, standard clockwise affine vectors
$v_1=(1,2,-4,0)^t$, $v_2=(-1,-3,3,0)^t$, and untwisted cusp collar. Thus $p=12l_0-4l_1-3l_2=-1$. We write $\mathfrak B=\PP(3,4)$, $\rho:\mathfrak B\to B$ for its coarse map, and $\bar f:X\to\mathfrak B$ for the root factorization. Let $\mathcal T=T_X(-\log W)$ and let $\nu:\mathrm dP_6\to X$ normalize $W$. The root generator is fixed by
\[
 M_n=\bar f^*\OO_{\mathfrak B}(n)
     =\OO_X(n(S_1-S_2)),\quad n\in\Z.
\]
The exponential isomorphism $H^1(X,\OO_X)\to\Pic(X)$ identifies $\Pic(X)$ with $\C$. We scale its coordinate so that $M_z=\exp(z\gamma^{0,1})$ and $M_1=\OO_X(S_1-S_2)$. The sign of this equality is fixed by the clockwise transgression $12\gamma\mapsto-p=1$. This is checked by evaluating the two finite lifts on $v_1,v_2$: their integral periods are $4l_1,3l_2$, while the cusp contributes $-12l_0$, whose sum is $-p$. The exponential sequence therefore gives $\exp(12\gamma^{0,1})=f^*\OO_B(1)$. Its injectivity and the root relation give the stated $M_1$ equality.

\begin{statement}{C5b}{Proved}{Integral twists and the normalization-map exception}
For $n\in\Z$ define
\[
 r_1(n)\equiv n\pmod3,\quad 0\le r_1(n)<3,
 \qquad r_2(n)\equiv3n\pmod4,\quad0\le r_2(n)<4,
\]
\begin{gather*}
 d(n)=\frac{n-4r_1(n)-3r_2(n)}{12},\qquad
 \epsilon_n=\begin{cases}0&3\mid n,\\1&3\nmid n.\end{cases}
\\
 u(n)=\max\{d(n)+1,0\},\qquad
 v(n)=\max\{-d(n)-1,0\}.
\end{gather*}
The coarse direct images of $T_X\otimes M_n$ are
\begin{align*}
 f_*(T_X\otimes M_n)&=\OO(d(n))\oplus\OO(d(n-1)),\\
 R^1f_*(T_X\otimes M_n)&=\OO(d(n))\oplus\OO(d(n-1))\oplus\OO(d(n+2)),\\
 R^2f_*(T_X\otimes M_n)&=\OO(d(n+6))\oplus\C_{p_1}^{\epsilon_n}.\tag{C.5}
\end{align*}
Their global cohomology dimensions are
\begin{align*}
 h^0&=u(n)+u(n-1),\\
 h^1&=v(n)+v(n-1)+u(n)+u(n-1)+u(n+2),\\
 h^2&=v(n)+v(n-1)+v(n+2)+u(n+6)+\epsilon_n,\\
 h^3&=v(n+6).                                             \tag{C.6}
\end{align*}
The natural map $H^2(X,T_X\otimes M_n)\to H^2(\mathrm dP_6,K_{\mathrm dP_6})=\C$ has rank one exactly for
\[
 n\in\{-6,-3,-2\}\cup\Z_{\ge0};                         \tag{C.7}
\]
it is zero for all other integers. In particular $M_{-1}$ has tangent cohomology $(0,0,1,0)$ and a zero map to the normalization target.
\end{statement}

\begin{proof}
We first check the precise input from R2--R3. The root charts at $p_1,p_2$ are the finite etale torus-family covers used in the construction; finite-group invariants over $\C$ are exact. At the cusp the local equations are $xy=q$ and $xyz=q$, and the vertical logarithmic fields are generated by $x\partial_x-y\partial_y$ and the remaining relative direction. Their affine transition derivatives depend only on the base. The invariant vertical field gives a trivial subline of the vertical bundle $E$. The quotient has twelfth power $\OO(-p_0)$: its automorphy is that of the reciprocal elliptic discriminant, whose only divisor is a simple pole at the cusp. As $\Pic(\mathfrak B)=\Z$, the quotient is $\OO(-1)$; $H^1(\mathfrak B,\OO(1))=0$ splits the extension. Thus $E=\OO\oplus\OO(-1)$.

The normalization complex of $W$ has global scalar complex
\[
 \C\xrightarrow{0}\C^3
 \xrightarrow{(a,b,c)\mapsto(a-b+c,a-b+c)}\C^2.
\]
The rational normalization and the three rational double curves have no positive-degree structure-sheaf cohomology. The displayed complex therefore gives $(1,2,1)$ for the fibre cohomology and its exterior cup product; the two oriented triangular faces give the nondegenerate degree-one pairing. The cyclic-cover fibres have the same dimensions. Proper flat base change on these charts gives locally free $G=R^1\bar f_*\OO_X$, with $R^2\bar f_*\OO_X=\det G$. Relative duality and the determinant of the logarithmic sequence give $\det G=\OO(-1)$. The integral class $\gamma$ maps to a nowhere-zero section of $G$: on a marked smooth torus the two identity columns of the period matrix prevent it from belonging to the span of the two holomorphic period rows, and at the cusp it is a nonzero cycle in the displayed complex. The same splitting argument yields
\[
E=G=\OO\oplus\OO(-1),\quad R^2\bar f_*\OO_X=\OO(-1),
\quad0\to\bar f^*E\to\mathcal T\to\bar f^*\OO(-5)\to0. \tag{C.8}
\]
These are the actual root-chart and cusp comparisons needed from the earlier returns; no coarse multiple fibre was replaced by its reduction in this argument.

Choose root sections $x$ of weight $3$, vanishing at $p_2$, and $y$ of weight $4$, vanishing at $p_1$. The Kodaira--Spencer morphism $\OO(-5)\to E\otimes G$ has, by its weights, the form
\[
 \kappa=\begin{pmatrix}0&a y\\b y&c_0x\end{pmatrix}.
\]
Its first coefficient vanishes because $H^0(\OO(5))=0$. The logarithmic period derivative at the cusp is invertible, giving $ab\ne0$. At $p_1$, the elliptic coordinate $\tau-\rho$ has a simple zero on the root chart, so its period derivative is nonzero, giving $c_0\ne0$. Rescaling the four line frames makes the three nonzero coefficients one. Cup product with this extension class gives
\[
\delta_0=(0,y,y,x)^t,\qquad
\delta_1=\begin{pmatrix}-y&0\\-x&y\end{pmatrix}:
 \OO(-5)\oplus\OO(-6)\to\OO(-1)\oplus\OO(-2).
\]
The first is a subbundle inclusion since $x,y$ have no common zero, and its cokernel is $\OO\oplus\OO(-1)\oplus\OO(2)$. The second is injective with determinant $-y^2$; its cokernel is
$Q=\OO(-1)|_{2p_1}$, a length-two module on the order-three chart. The direct-image sequence of (C.8) yields
\[
 R^0\bar f_*\mathcal T=E,\quad
 R^1\bar f_*\mathcal T=\OO\oplus\OO(-1)\oplus\OO(2),\quad
 R^2\bar f_*\mathcal T=\OO(-6)\oplus Q.
\]
The last splitting follows from $H^1(Q(6))=0$; it is not designated canonical.

Statement C5a gives the quotient $\nu_*K_{\mathrm dP_6}$. The local cusp volume form trivializes $\nu^*K_X$, and every $\nu^*M_z$ is trivial because $H^2(X,\Z)=0$ and $\Pic^0(\mathrm dP_6)=0$. Hence this quotient has cohomology only in degree two, where it is $\C$. At $p_0$, Serre duality and the conormal computation in Part A give $h^2(W,T_X|_W)=1$. Top-degree base change identifies this dimension with the residue dimension of $R^2\bar f_*T_X$. That module has generic rank one, so Nakayama's lemma and the structure theorem over $\C\{q\}$ make it free of rank one at $p_0$. Consequently its line summand is the elementary extension of $\OO(-6)$ by $\C_{p_0}$, namely $\OO(6)$, with no new cusp torsion. We have
\[
 R^0\bar f_*T_X=E,\quad
 R^1\bar f_*T_X=\OO\oplus\OO(-1)\oplus\OO(2),\quad
 R^2\bar f_*T_X=\OO(6)\oplus Q.                       \tag{C.9}
\]

For a root line, removing its two stabilizer characters gives
$\rho_*\OO(n)=\OO_B(d(n))$. Averaging the finite stabilizers is exact, so there are no higher direct images of $\rho$. Its sections are equivalently monomials $x^ay^b$ with $3a+4b=n$, and Serre duality gives $h^1\OO(n)=h^0\OO(-n-7)$. At $p_1$, the invariant exponent of $Q(n)$ is $n-1$ modulo three. Its two possible coefficients have orders zero and one, so its coarse length is $\epsilon_n$ and it is killed by the coarse parameter. Apply projection formula to (C.9) and then $\rho_*$. This proves (C.5), including the support and length of every torsion summand. Leray over $B$ has columns only zero and one, proving (C.6).

The normalization map on $R^2$ is the quotient in
\[
0\longrightarrow\OO(n-6)\longrightarrow\OO(n+6)
 \longrightarrow\C_{p_0}\longrightarrow0.
\]
Its global-section map is nonzero precisely when a section of $\OO(n+6)$ is nonzero at $p_0$. Both $x,y$ are nonzero at that point, so this occurs precisely when $n+6$ lies in the additive semigroup generated by $3,4$. That semigroup is $\{0,3,4,6,7,8,\ldots\}$, giving (C.7). The $H^1(R^1)$ part maps to zero in the cohomology of a sheaf supported at one point, by the natural Leray filtration. At $n=-1$, (C.5) has only one nonzero cohomological contribution, the reduced $\C_{p_1}$ summand in degree two, while $H^3(X,\mathcal T\otimes M_{-1})=H^1(\OO(-7))=\C$. The boundary from the normalization target is an isomorphism. This proves the stated counterexample to a surjection for every twist.
\end{proof}

\begin{statement}{C5c}{Proved}{Every nonintegral twist and its resonance lengths}
Retain the marked functions $\tau,\mu,\beta_c$ of the prompt and put $\rho=\tau(z_1)$. For $z\notin\Z$, let $\mathcal R_z$ be the set of regular base points at which integers $(n,p,q,r)$ satisfy
\[
 q\tau+r\mu=p,\qquad z=n-6q\mu-r\beta_c.               \tag{C.10}
\]
Let
\[
 \Lambda_1=\Z+(3\beta_c(z_1)+4-2\rho)\Z,\quad
 \Lambda_2=\Z+(2\beta_c(z_2)+3-3i)\Z.
\]
Set $\ell_1(z)=1+\epsilon_n$ for $z=n+k(3\beta_c(z_1)+4-2\rho)$ with $k\ne0$, and zero off $\Lambda_1$. Set $\ell_2(z)=1$ for $z=n+k(2\beta_c(z_2)+3-3i)$ with $k\ne0$ and $n\equiv1,2\pmod4$, and zero otherwise. Then $\mathcal R_z$ is finite. With $N_z=|\mathcal R_z|$ and $J_z=2N_z+\ell_1(z)+\ell_2(z)$, we have
\[
 (h^0,h^1,h^2,h^3)(X,T_X\otimes M_z)=(0,J_z,1+J_z,0),  \tag{C.11}
\]
and, writing $\mathfrak m_t$ for the maximal ideal of $\OO_{B,t}$,
\begin{align*}
R^1f_*(T_X\otimes M_z)&=
 \bigoplus_{t\in\mathcal R_z}\C_t^2
 \oplus\C_{p_1}^{\ell_1(z)}\oplus\C_{p_2}^{\ell_2(z)},\\
R^2f_*(T_X\otimes M_z)&=
 \C_{p_0}\oplus\bigoplus_{t\in\mathcal R_z}\OO_{B,t}/\mathfrak m_t^2
 \oplus\C_{p_1}^{\ell_1(z)}\oplus\C_{p_2}^{\ell_2(z)}.\tag{C.12}
\end{align*}
The natural map to $H^2(\mathrm dP_6,K_{\mathrm dP_6})$ is surjective for all these nonintegral twists.
\end{statement}

\begin{proof}
On a marked smooth torus, the exponential sequence presents $\Pic^0$ as constant complex lattice cocycles modulo the integral cocycles and the two holomorphic period rows. The equation for $z\gamma$ to belong to this sum is
\[
 z\gamma=(n,p,q,r)+a(6\mu,\tau,1,0)+b(\beta_c,\mu,0,1).
\]
The last two coordinates give $a=-q,b=-r$, and the first two give (C.10). If $r=0$, positivity of $\operatorname{Im}\tau$ forces $p=q=0$ and $z=n$. Otherwise
\[
 \operatorname{Im}z=-r\left(\operatorname{Im}\beta_c-
           6(\operatorname{Im}\mu)^2/\operatorname{Im}\tau\right)\ne0.
\]
In particular real nonintegral twists have no regular resonance. The fixed values $\mu(z_1)=(2-\rho)/3$ and $\mu(z_2)=(1-i)/2$ give the two displayed lattices by solving the first equation. Their nonreal generators make each representation $n+k\lambda_j$ unique.

We give a coherent local complex to calculate every length. On a smooth simply connected torus-family disc use normalized period coordinates $\zeta=Z(s)x+y$, $x,y\in\R^2/\Z^2$, and translation-invariant forms $G_j=(dx_j)^{0,1}$. The invariant horizontal lift is $h=\partial_s+\sum_i(Z'(s)x)_i\partial_{\zeta_i}$. Since $dx_j$ is closed, $\bar\partial G_j=0$ and
$\bar\partial h=\sum Z'_{ij}e_iG_j$. After subtracting the integer cocycle of a resonance, the line multiplier has invariant Dolbeault form $a_1G_1+a_2G_2$, up to the common nonzero factor $2\pi i$. Thus the invariant relative complex has terms
\[
 C^j=(V\oplus\OO h)\otimes\Lambda^jG,\qquad j=0,1,2,
\]
\[
 d(v\eta)=v(a\wedge\eta),\qquad
 d(h\eta)=h(a\wedge\eta)+\kappa\wedge\eta,             \tag{C.13}
\]
where
\[
 a=(z-n+6q\mu+r\beta_c,-p+q\tau+r\mu),\qquad
 \kappa=\begin{pmatrix}6\mu'&\tau'\\\beta_c'&\mu'\end{pmatrix}.
\]
This complex computes the derived direct image. To verify the comparison, decompose relative smooth forms into Fourier characters under real torus translations. The zero character gives (C.13). For every other character the relative Dolbeault differential is exterior multiplication by a nonzero covector; contraction by its normalized dual is a homotopy. On a sufficiently small closed subdisc those covectors are bounded away from zero at bounded frequency and grow linearly at large frequency. The homotopies preserve smooth Fourier series and all their derivatives. The vertical-horizontal filtration of the tangent extension, followed by the filtration in the base antiholomorphic differential, gives a finite filtration with acyclic nonzero-character quotients. Hence their full complexes are acyclic. The zero-character formulas have holomorphic coefficients, so the remaining base Dolbeault resolution is precisely the resolution of (C.13). A finite cyclic stabilizer normalizes translations, so this construction and its averaging projection are equivariant on each root chart.

At a regular resonance, $a'=\kappa^t(q,r)^t$ and $(q,r)\ne(0,0)$. The period determinant calculation in Part E shows that $\kappa$ is invertible at every regular point. Thus $a$ has a simple zero; choosing its direction as the second cocycle frame and changing the vertical frame gives
\[
 a=(0,s),\qquad\kappa=\begin{pmatrix}0&1\\1&0\end{pmatrix}.
\]
In degree one the two generators $e_1G_2,e_2G_2$ are killed by $s$. In degree two the relations eliminate the two vertical generators and leave $hG_1G_2$ killed by $s^2$. This gives the regular modules in (C.12).

At the order-three point put $t=s^3$, $\nu=\mu-(2-\tau)/3$ and $b=\beta_c+4\mu+(2/3)\tau-b(0)$. The order-three characters give $\nu=\nu_2s^2+O(s^3)$ and $b=b_3s^3+O(s^4)$. We check the joint coefficient that controls the module. With elliptic scale $f(s)=Fs^2+O(s^3)$, the Weierstrass equation gives
\[
 \nu_2^2=9F^2/512,\quad \tau_1=i\pi F^2/2,\quad
 6\mu'^2-\tau'\beta_c'=(3F^2/4)s^2+O(s^3).
\]
The first equality follows from the exact tripling coordinate $x(3P)(0)=512/81$ and $x=f^2\wp$: $3\nu$ approaches zero, so $x(3P)(0)=F^2/(9\nu_2^2)$. The determinant equality follows from the identity in Part E and $f^4g_2=12t^3(t-1)$, $f^6g_3=-8t^4(t-1)^2$. Comparing coefficients in $6\nu'^2-\tau'b'$ gives $b_3=7i/(32\pi)$. For a nonintegral resonance the Picard section is $3k(\nu,b)$ in the corresponding frames. Eliminating the unit $\tau'$ leaves transverse coefficient
\[
 3k(b-6\nu\nu'/\tau')
   =\frac{123ik}{64\pi}s^3+O(s^4),\qquad k\ne0.        \tag{C.14}
\]
It is nonzero. Equivariant holomorphic frame operations therefore reduce (C.13) to
\[
 a=(0,s^2),\qquad\kappa=\begin{pmatrix}0&s\\s&1\end{pmatrix}.
\]
Its degree-one and degree-two modules are both $\C\{s\}/(s^2)\oplus\C\{s\}/(s^3)$. Explicit degree-one cycles are $u=e_1G_2$, $v=e_2G_1+shG_2$; the third cycle is removed by the boundary $su+e_2G_2+sv$. The remaining relations are $s^2u=s^3v=0$. In degree two elimination of $e_2G_1G_2$ leaves the generators $e_1G_1G_2,hG_1G_2$ with annihilators $s^2,s^3$.

We retain the actual affine characters. The integral cocycle at the order-three resonance is $(n,-2k,-k,-3k)$; evaluating on $v_1$ gives $n$. At the order-four resonance it is $(n,-k,-k,-2k)$; evaluation on $v_2$ gives $-n\equiv3n\pmod4$. The invariant exponents at the order-three point are $n+2,n+2$ for the two degree-one generators and $n+2,n$ in degree two. A length-three summand has one invariant coefficient in every character; a length-two summand has one except for residue two. Both coarse lengths are therefore $1+\epsilon_n$, and multiplication by $t=s^3$ kills them. At the order-four point the regular normal form applies. The two length-one degree-one generators have invariant exponents $3n+1,3n+2$, and the length-two degree-two generator exponent $3n+2$. Their invariant length is one exactly for $n\equiv1,2\pmod4$, and the coarse parameter $s^4$ kills them. This proves all noncusp summands.

At the cusp, the restriction of $M_z$ has normalization pullback trivial and dual-complex character $(a,1)$, $a=e^{2\pi iz}$. Its normalization complex is
\[
 \C\xrightarrow{(a-1,0)^t}\C^2
 \xrightarrow{(0,a-1)}\C.
\]
For $z\notin\Z$ it is acyclic. Semicontinuity and base change make $R\bar f_*M_z$ vanish on a cusp neighbourhood. The logarithmic sequence (C.8) then gives $R\bar f_*(\mathcal T\otimes M_z)=0$ there. Sequence (C.3) contributes exactly the reduced skyscraper $\C_{p_0}$ in degree two for the ordinary tangent bundle. Away from resonances a nontrivial topologically trivial line on a torus has zero cohomology: its unitary Fourier presentation has no zero eigenmode; the tangent extension by trivial bundles retains that acyclicity. The direct-image support is consequently finite, since it is a proper analytic subset of a compact curve. Degree zero is zero because a direct image of a bundle is torsion-free and has generic rank zero. These arguments prove (C.12), and Leray proves (C.11). The logarithmic degree-two direct image is a torsion sheaf for these twists, so its first cohomology vanishes. Equation (C.4) proves the final surjectivity assertion.
\end{proof}

Statements C5b--C5c reproduce, with the displayed comparisons and independent exact module computations, the groups and stalk lengths claimed in S6\_01 D.1--D.4. The extra conclusion (C.7) corrects the all-twist interpretation of Lead 8: the target is always one-dimensional, but its global map is sometimes zero. All formulas give $\chi(T_X\otimes M_z)=1$. The generic tuple is $(0,0,1,0)$; it directly contradicts the degree-two vanishing in CDP20 Theorem 2.2(a), while its Euler characteristic and $c_3(X)=2$ contradict parts (b) and (c). A holomorphic universal Picard family is defined by $\bar\partial+z\alpha$ for any global closed Dolbeault representative $\alpha$ of the fixed generator. Proper coherent base change makes every cohomology threshold analytic. Real nonintegral parameters attain the generic tuple, so every higher threshold is a proper locally discrete analytic subset of the Picard line. This establishes the asserted jump-locus interpretation as well as the individual dimensions.
